% !TEX program = xelatex
% ============================================================
%  مخطط تقني — Schematic Sheet Template
%  قالب لورقة مخطط تقني (Schematic Sheet) بإطار حدود وكتلة عنوان
%  وخلفية شبكية لرسم المخططات، مع جدول تاريخ المراجعات وكتلة وسيلة الإيضاح.
%  Compile with: xelatex main.tex (run twice)
% ============================================================
%  Features:
%    - Border frame + title block (project, sheet, drawn/checked/approved by)
%    - TikZ grid background for drawing area
%    - Legend block for symbols
%    - Revision history table
% ============================================================

\documentclass[11pt,a4paper,landscape]{article}

\usepackage[a4paper,landscape,margin=1.2cm,top=1.2cm,bottom=1.2cm,headheight=14pt]{geometry}
\usepackage{graphicx}
\usepackage{array}
\usepackage{booktabs}
\usepackage{tabularx}
\usepackage{multirow}
\usepackage{enumitem}
\usepackage{float}
\usepackage{caption}
\usepackage{titlesec}
\usepackage{fancyhdr}
\usepackage{setspace}
\usepackage[table]{xcolor}
\usepackage{amsmath}
\usepackage{amssymb}
\usepackage{tikz}
\usetikzlibrary{positioning,calc,shapes.geometric,arrows.meta,decorations.pathreplacing}

\usepackage{bidi}
\usepackage[no-math]{fontspec}
\usepackage[quiet]{polyglossia}

\setmainfont[Scale=1]{Amiri-Regular.ttf}
\newfontfamily\arabicfont[Script=Arabic,Scale=0.95]{Amiri-Regular.ttf}
\newfontfamily\englishfont[Scale=0.95]{TeX Gyre Termes}

\setdefaultlanguage[calendar=gregorian,hijricorrection=1]{arabic}
\setotherlanguage[variant=british]{english}

\usepackage[breaklinks,hidelinks]{hyperref}
\hypersetup{pdftitle={مخطط تقني},pdfauthor={اسم الشركة}}

\addto\captionsarabic{%
  \renewcommand{\contentsname}{الفهرس}%
  \renewcommand{\figurename}{الشكل}%
  \renewcommand{\tablename}{الجدول}%
}

\titleformat{\section}{\large\bfseries}{\thesection}{0.6em}{}
\titlespacing*{\section}{0pt}{0.8ex plus 0.2ex}{0.5ex plus 0.2ex}

\pagestyle{fancy}
\fancyhf{}
\fancyhead[R]{\small\itshape مخطط تقني (Schematic Sheet)}
\fancyhead[L]{\small اسم المشروع}
\fancyfoot[C]{\small\thepage}
\renewcommand{\headrulewidth}{0.3pt}

\setlist[itemize]{topsep=0.2em, itemsep=0.1em, parsep=0pt, leftmargin=1.4em}

\newcolumntype{L}[1]{>{\raggedright\arraybackslash}p{#1}}
\newcolumntype{C}[1]{>{\centering\arraybackslash}p{#1}}
\newcolumntype{R}[1]{>{\raggedleft\arraybackslash}p{#1}}

\definecolor{frameColor}{HTML}{1A3C6B}
\definecolor{gridColor}{HTML}{D9D9D9}
\definecolor{titleBg}{HTML}{1A3C6B}
\definecolor{legendFill}{HTML}{F0F4F8}
\definecolor{accentColor}{HTML}{2E6DA4}

\setstretch{1.2}

\begin{document}

% ============================================================
% SCHEMATIC SHEET (border + grid + title block)
% ============================================================
% The whole sheet is drawn as a single TikZ canvas with an embedded
% grid drawing area, legend block, and title block.
% ============================================================
\thispagestyle{empty}

\begin{tikzpicture}[remember picture, overlay]
% ---- Outer border frame ----
\draw[line width=1.2pt, frameColor]
  ($(current page.north west)+(1.0cm,-1.0cm)$)
  rectangle
  ($(current page.south east)+(-1.0cm,1.0cm)$);

% Inner thin frame
\draw[line width=0.3pt, frameColor]
  ($(current page.north west)+(1.15cm,-1.15cm)$)
  rectangle
  ($(current page.south east)+(-1.15cm,1.15cm)$);

% ---- Drawing area with grid ----
\coordinate (drawTL) at ($(current page.north west)+(1.4cm,-1.4cm)$);
\coordinate (drawBR) at ($(current page.south east)+(-1.4cm,4.6cm)$);

\draw[line width=0.4pt, gridColor] (drawTL) grid[step=0.5cm] (drawBR);
\draw[line width=0.6pt, gray!60]   (drawTL) grid[step=2cm]   (drawBR);
\draw[line width=0.8pt, frameColor] (drawTL) rectangle (drawBR);

% Drawing area label
\node[anchor=north west, font=\small\color{frameColor}]
  at ($(drawTL)+(0.2cm,-0.2cm)$) {منطقة الرسم (\LR{Drawing Area})};

% ============================================================
% LEGEND BLOCK (bottom-left of drawing area)
% ============================================================
\coordinate (legTL) at ($(current page.south west)+(1.4cm,1.4cm)$);
\coordinate (legBR) at ($(legTL)+(7cm,3cm)$);

\fill[legendFill] (legTL) rectangle (legBR);
\draw[line width=0.6pt, frameColor] (legTL) rectangle (legBR);

% Legend header bar
\fill[titleBg] ($(legTL)+(0,2.6cm)$) rectangle ($(legTL)+(7cm,3cm)$);
\node[anchor=west, text=white, font=\small\bfseries]
  at ($(legTL)+(0.2cm,2.8cm)$) {وسيلة الإيضاح (\LR{Legend})};

% Legend entries
\node[anchor=west, font=\footnotesize] at ($(legTL)+(0.3cm,2.2cm)$)
  {\tikz\draw[line width=1pt, accentColor] (0,0) -- (0.8,0); \quad خط إشارة (\LR{Signal})};
\node[anchor=west, font=\footnotesize] at ($(legTL)+(0.3cm,1.7cm)$)
  {\tikz\draw[line width=1pt, dashed, accentColor] (0,0) -- (0.8,0); \quad خط تحكم (\LR{Control})};
\node[anchor=west, font=\footnotesize] at ($(legTL)+(0.3cm,1.2cm)$)
  {\tikz\draw[line width=1.4pt, frameColor] (0,0) -- (0.8,0); \quad خط طاقة (\LR{Power})};
\node[anchor=west, font=\footnotesize] at ($(legTL)+(0.3cm,0.7cm)$)
  {\tikz\fill[accentColor] (0,0) circle (2pt); \quad عقدة (\LR{Node})};
\node[anchor=west, font=\footnotesize] at ($(legTL)+(0.3cm,0.3cm)$)
  {\tikz\draw[line width=0.8pt] (0,0) rectangle (0.4,0.3); \quad مكون (\LR{Component})};

% ============================================================
% TITLE BLOCK (bottom-right)
% ============================================================
\coordinate (tbTL) at ($(current page.south east)+(-9.4cm,1.4cm)$);
\coordinate (tbBR) at ($(current page.south east)+(-1.4cm,4.6cm)$);

\draw[line width=0.8pt, frameColor] (tbTL) rectangle (tbBR);

% Title block header bar
\fill[titleBg] ($(tbTL)+(0,2.9cm)$) rectangle ($(tbBR)+(0,3.2cm)$);
\node[anchor=west, text=white, font=\small\bfseries]
  at ($(tbTL)+(0.2cm,3.05cm)$) {كتلة العنوان (\LR{Title Block})};

% Internal grid lines for title block
\draw[line width=0.3pt, frameColor]
  ($(tbTL)+(0,2.9cm)$)  -- ($(tbBR)+(0,2.9cm)$);  % under header
\draw[line width=0.3pt, frameColor]
  ($(tbTL)+(0,2.2cm)$)  -- ($(tbBR)+(0,2.2cm)$);
\draw[line width=0.3pt, frameColor]
  ($(tbTL)+(0,1.5cm)$)  -- ($(tbBR)+(0,1.5cm)$);
\draw[line width=0.3pt, frameColor]
  ($(tbTL)+(0,0.8cm)$)  -- ($(tbBR)+(0,0.8cm)$);
\draw[line width=0.3pt, frameColor]
  ($(tbTL)+(3.0cm,0)$)  -- ($(tbTL)+(3.0cm,2.9cm)$);
\draw[line width=0.3pt, frameColor]
  ($(tbTL)+(6.0cm,0)$)  -- ($(tbTL)+(6.0cm,2.9cm)$);

% Row 1: Project name (full width)
\node[anchor=west, font=\footnotesize\bfseries] at ($(tbTL)+(0.15cm,2.55cm)$) {اسم المشروع:};
\node[anchor=west, font=\footnotesize] at ($(tbTL)+(2.6cm,2.55cm)$) {---};

% Row 2: Sheet title | Sheet number
\node[anchor=west, font=\footnotesize\bfseries] at ($(tbTL)+(0.15cm,1.85cm)$) {عنوان الورقة:};
\node[anchor=west, font=\footnotesize] at ($(tbTL)+(2.6cm,1.85cm)$) {---};
\node[anchor=west, font=\footnotesize\bfseries] at ($(tbTL)+(3.15cm,1.85cm)$) {رقم الورقة:};
\node[anchor=west, font=\footnotesize] at ($(tbTL)+(5.6cm,1.85cm)$) {\LR{1 / 1}};

% Row 3: Drawn by | Checked by
\node[anchor=west, font=\footnotesize\bfseries] at ($(tbTL)+(0.15cm,1.15cm)$) {رسم بواسطة:};
\node[anchor=west, font=\footnotesize] at ($(tbTL)+(2.6cm,1.15cm)$) {---};
\node[anchor=west, font=\footnotesize\bfseries] at ($(tbTL)+(3.15cm,1.15cm)$) {مراجعة بواسطة:};
\node[anchor=west, font=\footnotesize] at ($(tbTL)+(5.6cm,1.15cm)$) {---};

% Row 4: Approved by | Date
\node[anchor=west, font=\footnotesize\bfseries] at ($(tbTL)+(0.15cm,0.45cm)$) {اعتماد بواسطة:};
\node[anchor=west, font=\footnotesize] at ($(tbTL)+(2.6cm,0.45cm)$) {---};
\node[anchor=west, font=\footnotesize\bfseries] at ($(tbTL)+(3.15cm,0.45cm)$) {التاريخ:};
\node[anchor=west, font=\footnotesize] at ($(tbTL)+(5.6cm,0.45cm)$) {---};

% Row 5: Company | Revision
\node[anchor=west, font=\footnotesize\bfseries] at ($(tbTL)+(0.15cm,-0.25cm)$) {اسم الشركة:};
\node[anchor=west, font=\footnotesize] at ($(tbTL)+(2.6cm,-0.25cm)$) {---};
\node[anchor=west, font=\footnotesize\bfseries] at ($(tbTL)+(3.15cm,-0.25cm)$) {المراجعة:};
\node[anchor=west, font=\footnotesize\bfseries] at ($(tbTL)+(5.6cm,-0.25cm)$) {\LR{Rev. A}};

\end{tikzpicture}

% ============================================================
% REVISION HISTORY (on a second page)
% ============================================================
\newpage

\begin{center}
{\LARGE\bfseries جدول تاريخ المراجعات (\LR{Revision History})}\\[0.5em]
{\large اسم المشروع --- مخطط تقني}\\[1em]
\end{center}

% TODO: املأ جدول تاريخ المراجعات.
\begin{table}[H]
\centering
\renewcommand{\arraystretch}{1.4}
\begin{tabularx}{\textwidth}{C{1.5cm} C{2.5cm} L{3cm} L{3cm} X}
\toprule
\textbf{المراجعة} & \textbf{التاريخ} & \textbf{المؤلف} & \textbf{المراجع} & \textbf{الوصف} \\
\midrule
\LR{A} & --- & --- & --- & الإصدار الأولي للمخطط. \\
\LR{B} & --- & --- & --- & % TODO: صف التغيير.
       وصف التغيير في المراجعة الثانية. \\
\LR{C} & --- & --- & --- & % TODO: صف التغيير.
       وصف التغيير في المراجعة الثالثة. \\
\bottomrule
\end{tabularx}
\end{table}

\vspace{1em}

\section{ملاحظات (\LR{Notes})}
% TODO: أضف ملاحظات فنية على المخطط.
\begin{itemize}
    \item جميع الأبعاد بالملليمتر (\LR{mm}) ما لم يُذكر خلاف ذلك.
    \item رموز المكونات وفق المعيار \LR{IEC 60617}.
    \item تُحدّث المراجعات في جدول تاريخ المراجعات أعلاه.
\end{itemize}

\end{document}
